% !TeX program = xelatex
% 题源均位于 27302模拟。下列简称及页码指对应 Markdown 标注的 PDF 页码：
% 880：27 李林 880 数二高数篇做题本（A4 紧凑版）。
% 500：李林基础 500 题数二题本（A4 版）。
% 六套卷：27 李林考研数学冲刺预测 6 套卷（数二）。
% 部分选择题调整了选项顺序。逐题出处见对应注释。
\documentclass[UTF8,a4paper,zihao=-4,fontset=fandol]{ctexart}
\usepackage[margin=20mm]{geometry}
\usepackage{amsmath,amssymb,enumitem,tabularx}
\setlength{\parindent}{0pt}
\setlength{\parskip}{0pt}
\setlength{\tabcolsep}{0pt}
\setlist[enumerate,1]{label=\arabic*.,leftmargin=2em}
\setlist[enumerate,2]{label=(\arabic*),leftmargin=2em,nosep}
\newcommand{\choices}[4]{%
  \par\vspace{0.15em}
  \begin{tabularx}{\linewidth}{@{}X@{\hspace{1em}}X@{}}
    \textup{(A)}\,#1 & \textup{(B)}\,#2 \\
    \textup{(C)}\,#3 & \textup{(D)}\,#4
  \end{tabularx}\par}
\newcommand{\blank}{\underline{\hspace{3em}}}
\raggedbottom

\begin{document}
\begin{center}
  {\LARGE\bfseries 27302 高等数学阶段模拟卷\par}
  \vspace{0.4em}
  {\small 《工科数学分析基础》第三版上册\quad 第一章至第二章第四节（微分中值定理及其应用）\par}
  \vspace{0.4em}
  {\small 满分 150 分\qquad 建议用时 180 分钟}
\end{center}

\section*{一、选择题}
第 1--10 题，每题 5 分，共 50 分。每题只有一个选项符合题意。

\begin{enumerate}[itemsep=0.4em,topsep=0.5em]
  % 880，第 2 页，第一章基础题，选择题 (3)。
  \item 设 $f(x)=\sqrt{1+x+x^2}-\sqrt{1-x+x^2}$，则（\quad）。
  \choices{$f(x)$ 为偶函数}{$f(x)$ 为奇函数}
          {$f(x)$ 为无界函数}{$\displaystyle\lim_{x\to-\infty} f(x)=1$}

  % 880，第 7 页，第一章综合题，选择题 (11)。
  \item 设数列 $\{x_n\}$ 单调递减，$\{y_n\}$ 单调递增，且
  $\displaystyle\lim_{n\to\infty}(x_n-y_n)=0$，则（\quad）。
  \choices{$\{x_n\}$ 与 $\{y_n\}$ 均收敛于同一极限}
          {$\{x_n\}$ 与 $\{y_n\}$ 均收敛于 $0$}
          {$\{x_n\}$ 收敛，$\{y_n\}$ 发散}
          {$\{x_n\}$ 与 $\{y_n\}$ 均发散}

  % 880，第 7 页，第一章综合题，选择题 (7)。
  \item 设 $f(x)=
    \begin{cases}
      (x+1)\arctan\dfrac{1}{x^2-1}, & x\ne\pm1,\\
      0, & x=\pm1,
    \end{cases}$
    则 $f(x)$（\quad）。
  \choices{在 $x=1$ 与 $x=-1$ 处均连续}
          {在 $x=1$ 与 $x=-1$ 处均间断}
          {在 $x=-1$ 处间断，在 $x=1$ 处连续}
          {在 $x=-1$ 处连续，在 $x=1$ 处间断}

  % 六套卷，第 16 页，第四套第 1 题。
  \item 设常数 $a,b$ 满足 $\displaystyle\lim_{x\to+\infty}(x\arctan x-ax-b)=0$，则（\quad）。
  \choices{$a=1,\ b=1$}{$a=\pi,\ b=-1$}
          {$a=\dfrac{\pi}{2},\ b=-1$}{$a=\dfrac{\pi}{2},\ b=1$}

  % 500，第 16 页，第一章概念与理解篇第 13 题。
  \item 设 $f$ 和 $\varphi$ 在 $\mathbb{R}$ 上有定义，$f$ 连续，
  $x=a$ 是 $\varphi$ 的跳跃间断点。下列说法错误的是（\quad）。
  \choices{$f\varphi$ 在 $x=a$ 处可能连续}
          {$f\varphi$ 在 $x=a$ 处可能有可去间断点}
          {$f\circ\varphi$ 在 $x=a$ 处可能连续}
          {$f\circ\varphi$ 在 $x=a$ 处可能有可去间断点}

  % 500，第 17 页，第二章计算与方法篇第 1 题。
  \item 设 $f(x)=x(x+1)\cdots(x+20)$，则（\quad）。
  \choices{$f'(-1)>0,\ f'(-2)>0$}{$f'(-1)>0,\ f'(-2)<0$}
          {$f'(-1)<0,\ f'(-2)<0$}{$f'(-1)<0,\ f'(-2)>0$}

  % 880，第 15 页，第二章基础题，选择题 (7)。
  \item 设 $y=f(x)$ 可导且 $f'(x_0)=\dfrac12$。当 $\Delta x\to0$ 时，
  $f$ 在 $x_0$ 处的微分 $\mathrm{d}y$ 是 $\Delta x$ 的（\quad）无穷小。
  \choices{同阶但不等价}{等价}{低阶}{高阶}
\end{enumerate}

\newpage
\begin{enumerate}[start=8,itemsep=0.4em,topsep=0pt]
  % 880，第 24 页，第二章综合题，选择题 (5)。
  \item 设 $f$ 在 $x=0$ 的某邻域内有定义。函数 $f(x)|\sin x|$ 在 $x=0$ 处可导的充要条件是（\quad）。
  \choices{$\displaystyle\lim_{x\to0}f(x)$ 存在}
          {$\displaystyle\lim_{x\to0}f(x)=f(0)$}
          {$f$ 在 $x=0$ 处可导}
          {$f$ 在 $0$ 处的两个单侧极限均存在且互为相反数}

  % 880，第 16 页，第二章基础题，选择题 (19)。
  \item 曲线 $y=x^3+ax^2+bx+c$ 在 $x=-2$ 处取得极值，
  且在点 $(1,0)$ 处与直线 $y=-3x+3$ 相切，则（\quad）。
  \choices{$a=1,\ b=-8,\ c=6$}{$a=-1,\ b=-8,\ c=-6$}
          {$a=1,\ b=8,\ c=-6$}{$a=-1,\ b=8,\ c=-6$}

  % 880，第 24 页，第二章综合题，选择题 (4)。
  \item 设 $f$ 在 $[0,1]$ 上二阶可导，$f(0)=f(1)$，且对任意 $x\in(0,1)$ 有
  $f''(x)\ne0$，则（\quad）。
  \choices{$f$ 在 $(0,1)$ 内至少有一个零点}{$f'$ 在 $(0,1)$ 内没有零点}
          {$f'$ 在 $(0,1)$ 内有唯一零点}{$f'$ 在 $(0,1)$ 内至少有两个零点}
\end{enumerate}

\section*{二、填空题}
第 11--16 题，每题 5 分，共 30 分。将答案填在题中横线上。

\begin{enumerate}[start=11,itemsep=0.6em]
  % 880，第 9 页，第一章综合题，填空题 (4)。
  \item 已知 $\displaystyle\lim_{x\to a}\frac{f(x)-b}{x-a}=A$，
  则 $\displaystyle\lim_{x\to a}\frac{\mathrm{e}^{f(x)}-\mathrm{e}^{b}}{x-a}=\blank$。

  % 六套卷，第 2 页，第一套第 11 题。
  \item 设 $f$ 在 $x=1$ 处连续，且 $\displaystyle\lim_{x\to1}\frac{f(x)-5x+2}{\ln x}=1$，
  则 $f'(1)=\blank$。

  % 500，第 26 页，第二章计算与方法篇第 28 题。
  \item 设
  $\begin{cases}
    x=f(\arctan t)-\pi,\\
    y=f(\mathrm{e}^{-2t}-1),
  \end{cases}$
  其中 $f$ 可导且 $f'(0)\ne0$，则 $\left.\dfrac{\mathrm{d}y}{\mathrm{d}x}\right|_{t=0}=\blank$。

  % 500，第 26 页，第二章计算与方法篇第 30 题。
  \item 设 $y=y(x)$（$0<y<\pi$）由方程 $y\sin x-\cos(x+y)=0$ 确定，
  则 $y''(0)=\blank$。

  % 500，第 29 页，第二章计算与方法篇第 38 题。
  \item 对 $x>-1$ 且 $x\ne0$，设 $\xi=\xi(x)$ 为 $f(t)=\ln(1+t)$
  在以 $0,x$ 为端点的区间上满足拉格朗日中值定理的中值点，
  则 $\displaystyle\lim_{x\to0}\frac{\xi(x)}{x}=\blank$。

  % 880，第 30 页，第二章综合题，填空题 (7)。
  \item 设 $f$ 在 $x_0$ 处有三阶连续导数，且
  $f'(x_0)=1$，$f''(x_0)=2$，$f'''(x_0)=3$。
  设 $x=g(y)$ 是 $y=f(x)$ 在 $x_0$ 附近的反函数，$y_0=f(x_0)$，则 $g'''(y_0)=\blank$。
\end{enumerate}

\newpage
\section*{三、解答题}
第 17 题 10 分，第 18--22 题每题 12 分，共 70 分。写出必要的演算步骤或证明过程。

\begin{enumerate}[start=17]
  % 880，第 5 页，第一章基础题，解答题 (2) 第 (V)、(VI) 问。
  \item \textbf{（10 分）} 求下列极限。
  \[
    \text{(1)}\quad\lim_{x\to0}\frac{\mathrm{e}^{\tan x}-\mathrm{e}^{x}}{x^3},
    \qquad
    \text{(2)}\quad\lim_{x\to0}\cot x\left(\frac{1}{\sin x}-\frac{1}{x}\right).
  \]
  \vspace{5.5cm}

  % 880，第 21 页，第二章基础题，解答题 (5)。
  \item \textbf{（12 分）} 设曲线由参数方程
  $x=t-\sin t,\ y=1-\cos t$ 确定。对 $t\notin2\pi\mathbb{Z}$，
  求 $\dfrac{\mathrm{d}y}{\mathrm{d}x}$ 和 $\dfrac{\mathrm{d}^2y}{\mathrm{d}x^2}$。
  \vspace{5.5cm}

  % 880，第 31 页，第二章综合题，解答题 (4)。
  \item \textbf{（12 分）} 设
  \[
    f(x)=
    \begin{cases}
      |x|^p\sin\dfrac{1}{x}, & x\ne0,\\
      0, & x=0,
    \end{cases}
    \qquad p\in\mathbb{R}.
  \]
  分别求使 $f$ 在 $x=0$ 处连续、可导，以及使 $f'$ 在 $x=0$ 处连续的 $p$ 的取值范围。
\end{enumerate}

\newpage
\begin{enumerate}[start=20]
  % 880，第 11 页，第一章综合题，解答题 (4)。
  \item \textbf{（12 分）} 设 $f$ 在 $[0,1]$ 上连续，且 $f(0)=f(1)$。证明下列结论。
  \begin{enumerate}
    \item 存在 $\xi\in(0,\tfrac12]$，使 $f(\xi)=f(\xi+\tfrac12)$。
    \item 对任意整数 $n\ge2$，存在 $\xi\in(0,1-\tfrac1n]$，使 $f(\xi)=f(\xi+\tfrac1n)$。
  \end{enumerate}
  \vspace{6cm}

  % 880，第 33 页，第二章综合题，解答题 (22)。
  \item \textbf{（12 分）} 设 $f$ 在 $[0,1]$ 上连续，在 $(0,1)$ 内可导，
  且 $f(0)=0$、$f(1)=1$。证明下列结论。
  \begin{enumerate}
    \item 存在 $0<\xi_1<\xi_2<1$，使 $f'(\xi_1)+f'(\xi_2)=2$。
    \item 存在 $\xi,\eta\in(0,1)$，使 $\eta f'(\xi)=f(\eta)f'(\eta)$。
  \end{enumerate}
  \vspace{6cm}

  % 880，第 31 页，第二章综合题，解答题 (9)。
  \item \textbf{（12 分）} 设 $a<b$，$f$ 在 $[a,b]$ 上可导。
  记 $f'_+(a)$ 和 $f'_-(b)$ 分别为 $f$ 在两个端点处的右导数和左导数。
  \begin{enumerate}
    \item 若 $f'_+(a)f'_-(b)<0$，证明存在 $\xi\in(a,b)$，使 $f'(\xi)=0$。
    \item 若 $f'_+(a)\ne f'_-(b)$，证明对严格介于二者之间的任意实数 $\mu$，
    均存在 $\xi\in(a,b)$，使 $f'(\xi)=\mu$。
  \end{enumerate}
\end{enumerate}
\end{document}
